\documentclass[10pt,conference,compsoc]{IEEEtran}

% ── Core packages ──────────────────────────────────────────────────
\usepackage[utf8]{inputenc}
\usepackage[T1]{fontenc}
\usepackage[margin=0.7in]{geometry}
\usepackage{titlesec}
\usepackage{amsmath,amssymb,amsthm,mathtools,bm}
\usepackage{graphicx,xcolor,hyperref}
\usepackage{booktabs,enumitem,caption,subcaption,array}
\usepackage{algorithm,algpseudocode}

\hypersetup{colorlinks=true,linkcolor=blue!70!black,
            citecolor=green!50!black,urlcolor=blue!60!black}
\setlist{nosep}

% ── Math macros ─────────────────────────────────────────────────────
\newcommand{\R}{\mathbb{R}} \newcommand{\E}{\mathbb{E}}
\newcommand{\Var}{\operatorname{Var}} \newcommand{\Cov}{\operatorname{Cov}}
\newcommand{\VR}{\operatorname{VR}}   \newcommand{\sign}{\operatorname{sign}}
\newcommand{\bw}{\bm{w}}  \newcommand{\br}{\bm{r}}
\newcommand{\MOM}{\textsc{mom}}  \newcommand{\REV}{\textsc{rev}}

% ── Title block (tight) ────────────────────────────────────────────
\title{\large \textbf{Project 1: Mean-Reversion vs.\ Momentum Across Horizons:\\
A Data-Source-Matched Comparison on the Dow Jones 30}\vspace{-0.5em}}

\author{
  \small E. Pan, L. Miniutti, A.\ Favara --- \textit{Massachusetts Institute of Technology
  (15.C51 Spring 2026), Proprietary Trading Track}\vspace{-5em}
}

\begin{document}

\titlespacing*{\section}{0pt}{1.2ex}{0.6ex}
\titlespacing*{\subsection}{0pt}{1.0ex}{0.5ex}
\titlespacing*{\subsubsection}{0pt}{0.8ex}{0.4ex}
\titleformat{\paragraph}[runin]{\normalfont\bfseries}{}{0pt}{}
\titlespacing{\paragraph}{0pt}{0.5ex}{1em}

\maketitle
\setlength{\parskip}{0pt}

%======================================================================
\section{Introduction}
%======================================================================

The classical equity anomalies literature disagrees about what individual
stocks do to their own past returns. Short-horizon studies find
\textbf{reversal}: a stock that moves sharply today partially reverses
tomorrow~\cite{lehmann1990}, over a week~\cite{lo1990}, or over a
month~\cite{jegadeesh1990}; De~Bondt--Thaler~\cite{dt1985} extend this to
three- and five-year windows. Intermediate-horizon studies find
\textbf{momentum}: three-to-twelve-month winners keep
winning~\cite{jt1993}. Very short (bar-level) horizons are dominated by
bid-ask bounce~\cite{roll1984}, the hardest regime to profit in. At the
daily boundary, Gao, Han, Li, and Zhou~\cite{gao2018} find positive
intraday-to-close autocorrelation on S\&P~500 futures.

These studies differ on universe, horizon, and data vintage. We try to run that comparison on the \textbf{Dow Jones 30} over
\textbf{2016-01-04 to 2025-09-30}, rebuilding each horizon on the
dataset that natively serves it. Namely:
\begin{quote}
\itshape At which horizons does a cross-sectional tercile long--short
on past returns produce significant \emph{positive} net Sharpe, and by
how much does each fail or succeed relative to its classical priors?
\end{quote}
We hope the sweep offers some discrimination between competing
mechanisms: different horizons correspond to different economic
stories in the literature, and running them side-by-side on a
single universe with a uniform portfolio construction lets us see
which stories survive in this sample and which do not. If every
horizon had failed, we would have reported a null result; if every
horizon had succeeded, we would have spent more time checking the
code before drawing conclusions. What we actually find --- one
horizon clearing significance, the others failing for distinct and
identifiable reasons (costs, large-cap liquidity, overlap
correction) --- is a modest but informative pattern that a single
horizon study could not have produced.

We apply a \textbf{single portfolio construction} (equal-weight,
dollar-neutral tercile long--short) at six horizons spanning nine orders
of magnitude in time: H1 30-minute bars, H2 intraday first-30 to
rest-of-day, and H3--H6 daily-to-semiannual. We run both directions (\MOM{}: long winners, short losers; \REV{}: reverse) and control the resulting family of 12 tests at $\alpha = 5\%$ with a Bonferroni threshold $t^{*}\!\approx 2.87$.

Five findings anchor the rest of this report:

\begin{enumerate}[leftmargin=*]
  \item \textbf{H2 REV is the one Bonferroni-significant strategy
    in the family:} net Sharpe $+2.53$, $t = +7.86$, max drawdown
    $-4.3\%$, comfortably above the family-wise threshold
    (Figure~\ref{fig:horizon}).
  \item \textbf{H2 REV has near-zero exposure to standard factors.}
    A daily Fama-French 5-factor plus momentum regression gives
    $\hat\alpha = +10.3\%$ annualised ($t = +7.82$) with every
    factor loading insignificant and $R^{2} = 0.003$
    (Table~\ref{tab:factors}). The return stream is largely
    unexplained by standard daily risk factors --- consistent with
    an intraday strategy that holds no overnight risk.
  \item \textbf{Every other horizon falls short of the family-wise
    threshold}, for reasons we can identify. H3 reversal is
    eaten by transaction costs at daily frequency; H5 one-month
    reversal does not appear in this DJ30 sample; H6
    Jegadeesh--Titman momentum has the right sign (net SR $+0.82$)
    but its overlap-adjusted $t$-stat (Newey--West, $q{=}6$)
    falls from $+2.43$ to $+1.37$, below Bonferroni.
  \item \textbf{H1 (30-minute bars) is dominated by transaction
    costs.} Gross Sharpe is $+0.05$; at 3~bps per side
    $\times$ $1.35$ bar-over-bar turnover $\times$ 2{,}772 bars/year,
    cost drag is $\sim$112\% annualised and net wealth decays to
    near zero. The result is consistent with the prior expectation
    that tercile rebalancing at 30-minute cadence is not
    cost-feasible on DJ30 large caps.
  \item \textbf{Time-series and cross-sectional measures of
    reversal are complementary rather than equivalent.} The
    Lo--MacKinlay variance ratio $\mathrm{VR}(q)$ --- a
    per-stock, time-series measure --- falls monotonically from
    $0.95$ at $q{=}2$d to $0.56$ at $q{=}252$d (median across 40
    tickers; Figure~\ref{fig:vr}), clearly showing reversal
    signatures in each stock's own return series. Our
    cross-sectional tercile long--short asks a different question
    and, as the headline shows, yields a profitable answer only at
    H2. The two views together give a more complete picture than
    either alone.
\end{enumerate}

\paragraph{Contribution.} The study offers a data-source-matched
horizon sweep on a fixed universe with a uniform methodology, a
finding for individual DJ30 names that diverges in sign from Gao et
al.'s~(2018) index-level intraday-momentum result, and a robustness
suite (cost sensitivity, Roll bid-ask-bounce diagnostic, VIX-regime
/ year-by-year stability, survivorship, look-ahead audit, FF5+Mom
factor regression) that supports the H2 REV finding within the
confines of this sample.

%======================================================================
\section{Data \& Universe}
\label{sec:data}
%======================================================================

\subsection{Three datasets, three horizon domains}

Each horizon is built on the dataset that natively serves it, so a
single research question is answered without cross-database splicing.

\paragraph{CRSP Daily.} The Center for Research in Security
Prices' Daily Stock File, 2016-01-04 to 2025-09-30. For every DJ30
ever-member we use \texttt{DlyRet} (total return including dividends),
\texttt{DlyPrcVol} (dollar volume) and prices (\texttt{abs(DlyClose)}
etc.\ to handle CRSP's negative-midquote convention). Serves horizons
H3--H6.

\paragraph{WRDS Intraday Indicators (iid\_ms).} A pre-aggregated
daily panel of microstructure statistics computed from
millisecond-stamped TAQ, one row per (ticker, date). We use
\texttt{OPrc}, \texttt{CPrc}, \texttt{mid\_after\_open}
(a $\sim$5-minute VWAP midprice immediately after the opening auction),
plus quoted spreads and quote-based intraday vols. Serves H2.

\paragraph{TAQ Millisecond.} Consolidated-trades files, one per
calendar year, 93~GB gzipped total across 2016--2025. Processed on
the MIT Engaging cluster to 30-minute OHLCV bars per (ticker,
date). Serves H1.

\subsection{Point-in-time universe}

The DJ30 changed members ten times in our window. Using today's
membership would introduce survivorship bias: NVDA, AMZN, SHW, HON,
CRM, and AMGN were all added after 2016. We built a point-in-time
membership file of 75{,}420 (date, ticker) pairs (2{,}514 trading
days $\times$ exactly 30 members each) from S\&P Dow Jones Indices
press releases. Every horizon panel is inner-joined against this
file before backtesting, and the 30-names assertion is enforced
modulo a short list of ticker-chain boundary days (DD$\to$DWDP,
DWDP$\to$DOW, UTX$\to$RTX).

A concrete illustration of why this matters: NVDA returned roughly
$40\times$ from 2016 to 2025 but only joined the DJ30 in November
2024. A non-PIT universe would give a momentum strategy a 4{,}000\%
tailwind on a name that was not actually in the index for 99\% of
the backtest window. The PIT filter removes this bias entirely; it
is the single data-engineering step with the biggest influence on
the final headline.

\subsection{Two data-quality fixes worth recording}

An original WRDS extract of 40 ever-member tickers was missing
Goldman Sachs (PERMNO~86868), which happens to be a DJ30 member on
\emph{every} trading day in the window. A supplementary WRDS pull of
PERMNO~86868 was concatenated into the CRSP and iid loaders; the
loaders use the supplement file if it exists, so a future full re-pull
is backward-compatible. Separately, at the UTC+Raytheon merger close
(Fri 2020-04-03), CRSP and iid both relabel PERMNO~17830 from UTX to
RTX immediately, but S\&P DJI kept UTX in the index through that day
and added RTX the following Monday. A single-cell ticker patch in the
loader aligns the two conventions; the underlying security is the same
merged entity.

%======================================================================
\section{Methodology}
\label{sec:method}
%======================================================================

\subsection{Signal and forward-return definitions}

At rebalance date $t$ and horizon $h$, let $U_t$ be the 30-member DJ30
cross-section as of date $t$. For each $i \in U_t$ we compute a
\textbf{signal} $s_{i,t}^{(h)}$ and a \textbf{forward return}
$r_{i,t+h}^{\text{fwd}}$. The methodology is uniform across horizons;
only the windowing differs.

\paragraph{Signal.} The past-horizon return,
\begin{equation}
  s_{i,t}^{(h)} = \prod_{k \in W_{\text{past}}(h,t)}
                  \bigl(1 + r_{i,k}\bigr) - 1,
  \label{eq:signal}
\end{equation}
where $W_{\text{past}}$ is the horizon-specific lookback window. For
H3 it is the single day $t$; for H5 the month ending at $t$; for H6
the 105 trading days ending 21 days before $t$ (Jegadeesh--Titman
``skip-1-month'' convention~\cite{jt1993}).

\paragraph{Forward.} The cumulative return to the next rebalance,
\begin{equation}
  r_{i,t+h}^{\text{fwd}} =
  \prod_{k \in W_{\text{fwd}}(h,t)}\bigl(1 + r_{i,k}\bigr) - 1.
  \label{eq:fwd}
\end{equation}
For non-overlapping horizons (H3, H4, H5), $W_{\text{fwd}}$ is the
trading days from $t{+}1$ to the next rebalance. For H6 it is the
fixed 126-day window $[t{+}1, t{+}126]$, overlapping with subsequent
rebalances (see \S\ref{sec:nw}). For H2 both windows are intraday
(same date), with $s^{(H2)} = \text{mid\_after\_open}/\text{OPrc} - 1$
and $r^{\text{fwd}} = \text{CPrc}/\text{mid\_after\_open} - 1$.

\subsection{Tercile long--short portfolio}

At each $t$, rank $U_t$ by $s^{(h)}_{i,t}$ with rank ties broken
stably. Let $q_{i,t} \in \{\text{lo},\text{mi},\text{hi}\}$ be tercile
membership. The \MOM{} weights are
\begin{equation}
  w^{\MOM}_{i,t} = \frac{1}{2\,|T_{q_{i,t}}|}\cdot
  \begin{cases}
    +1 & q_{i,t} = \text{hi}\\
    -1 & q_{i,t} = \text{lo}\\
    \phantom{-}0 & q_{i,t} = \text{mi},
  \end{cases}
  \label{eq:weights}
\end{equation}
where $|T_{q}|$ is the size of tercile $q$ on date $t$. The \REV{}
weights are the exact mirror: $w^{\REV} = -w^{\MOM}$. By construction,
$\sum_i w_{i,t} = 0$ (dollar-neutral) and $\sum_i |w_{i,t}| = 1$
(gross leverage one). On a 30-name universe terciles are $10/10/10$;
on short days (ticker transitions) they adjust to the
closest-to-equal split.

\paragraph{Why terciles on a 30-name universe.} Three choices
compete: terciles (top/bottom 33\%, 10 names per leg), quintiles
(top/bottom 20\%, 6 names per leg), deciles (top/bottom 10\%, 3
names per leg). The per-date noise penalty of smaller legs is
$\sqrt{10/|T_q|}$ --- about $1.29\times$ for quintiles,
$1.83\times$ for deciles. The signal gain is the increase in
mean-signal separation between extreme groups, which is of order
$\sigma_s / \sqrt{|T_q|}$ for approximately-normal signal
distributions. The ratio of gain to penalty favours terciles for
near-linear signal--return relationships (our case at most
horizons) and flips toward smaller legs only if the signal--return
curve is sharply convex at the extremes. Deciles are also
untradeable: three-name legs make the portfolio hostage to any
single-name idiosyncratic shock (an earnings day, a CEO
resignation).

\paragraph{Why rank-sort instead of signal-weighted.} Table~II
(panel integrity) shows every horizon's signal has positive excess
kurtosis (up to $+32.7$ on H2). A signal-weighted portfolio
(e.g., $w_i \propto s_i$) would concentrate its bets on the day's
most extreme movers --- the 99th-percentile news days. That makes
the strategy's Sharpe sensitive to a small number of tail events,
which we would rather avoid. The rank sort maps every name to a
tercile regardless of signal magnitude, which keeps per-day P\&L
dispersion tight and the effective sample size high.

\subsection{Cost model}

Two conventions. For horizons that hold positions across rebalances
(H3--H6), turnover is incremental:
\begin{equation}
  \tau_{t+h}^{\text{inc}} =
  \sum_{i \in U_t \cup U_{t+h}} \bigl|w_{i,t+h} - w_{i,t}\bigr|.
  \label{eq:turn}
\end{equation}
For the intraday horizon H2, positions open at $\sim$10:00 and fully
unwind at 16:00; between days we are flat. Turnover is then fixed at
$\tau^{\text{rt}} = 2\cdot\sum_i |w_{i,t}| = 2$ per rebalance. Net
P\&L in both cases is
\begin{equation}
  \Pi^{\text{net}}_{t+h} \;=\;
  \sum_i w_{i,t}\, r_{i,t+h}^{\text{fwd}}
  \;-\; c_h\cdot\tau_{t+h},
  \label{eq:net}
\end{equation}
with $c_h$ the per-side cost in decimal (i.e., $1.5~\text{bps} =
1.5\times10^{-4}$). Table~\ref{tab:horizons} lists per-horizon
$c_h$, set by typical DJ30 half-spread plus impact at the applicable
trading frequency.

\paragraph{Mental model for the two cost conventions.} $c_h$ is
what we pay to cross the spread (plus impact) one-way; a round-trip
(open + close) costs $2c_h$. H2's intraday convention always incurs
a full round-trip per day because positions go from flat to
target at $\sim$10:00 and back to flat at 16:00. H3--H6 carry
positions forward, so a rebalance only pays on the \emph{changed}
portion of the portfolio; if ranks are stable bar-over-bar, the
cost is small. This is why H6 MOM has the lowest cost drag per
rebalance (mean turnover $0.56$) and H1 the highest cost drag per
year (turnover $1.35$ per \emph{bar} $\times$ 2{,}772 bars/year
$=$ 3{,}744 of turnover per year, i.e.\ 30 full flips of the book
weekly).

\subsection{Newey--West for H6 overlap}
\label{sec:nw}

Consecutive H6 monthly rebalances have forward windows that overlap
by 5 months, inducing autocorrelation in the P\&L series at lags
1--5. Naive $t$-statistics are inflated by up to $\sqrt{1+2\cdot
5/6}\approx 1.63\times$. We compute both a naive $t$ and a
Newey--West HAC-adjusted $t$ with $q=6$ lags~\cite{nw1987}:
\begin{equation}
  \hat V_{\text{NW}}
  \;=\; \hat{\gamma}_0 + 2\sum_{k=1}^{q}\Bigl(1-\tfrac{k}{q+1}\Bigr)
         \hat{\gamma}_k,
  \label{eq:nw}
\end{equation}
where $\hat\gamma_k$ is the $k$-th-lag autocovariance of the
demeaned P\&L series. The Bonferroni threshold is applied to the
corrected $t$ for H6 and the naive $t$ elsewhere.

\paragraph{Why $q=6$.} Andrews~(1991) recommends
$q \propto T^{1/3}$ for HAC truncation lag; with $T \approx 100$
monthly observations that prescribes $q \approx 5$--$6$. The H6
forward window is specifically a 6-month hold, so $q=6$ also
covers the full known autocovariance structure. Appendix~B shows
the H6 corrected $t$ is stable at $1.35$--$1.53$ across
$q \in \{4, \ldots, 18\}$ --- the conclusion (H6 does not clear
Bonferroni) does not depend on the specific lag choice.

\subsection{Bonferroni multiple-testing control}

The study runs $n = 12$ tests ($6$ horizons $\times$ $2$ directions).
Controlling family-wise error rate at $\alpha = 5\%$,
\begin{equation}
  t^{*} \;=\; \Phi^{-1}\!\bigl(1 - \tfrac{\alpha/n}{2}\bigr)
  \;\approx\; 2.8653.
  \label{eq:bonf}
\end{equation}
A strategy's $t$-stat must exceed $+t^{*}$ on the correct side to be
declared a Winner; two-sided thresholds would include losing strategies
with very negative $t$ (e.g.\ H1 at $|t|\approx 91$), which is not
economically meaningful.

\paragraph{Why Bonferroni rather than Benjamini--Hochberg.} BH
controls the false-\emph{discovery} rate (the expected proportion
of false positives among rejected nulls), which is appropriate for
exploratory science where some false positives are tolerable. For
trading, a false positive is a losing strategy in live capital ---
we want control over the probability of making \emph{any} false
discovery at all, which is exactly the family-wise error rate that
Bonferroni controls. Bonferroni is also conservative: a strategy
that fails it but passes BH is a strategy whose significance
vanishes once you correct for the other eleven tests we ran,
which is not a strategy we want to recommend.

\subsection{Backtest engine in full}
\label{sec:algo}

Algorithm~\ref{alg:engine} collapses the methodology into the exact
procedure every horizon runs through. The function is pure-pandas with
three hot paths: per-date ranking, per-date P\&L summation, and a
wide-pivot turnover diff. It is the single primitive shared across all
twelve ($6 \times 2$) backtests in this report.

\begin{algorithm}[ht]
\small
\caption{\textsc{run\_horizon} --- unified tercile long--short}
\label{alg:engine}
\begin{algorithmic}[1]
\Require panel $\mathcal{P} = \{(i, t, s_{i,t}, r^{\text{fwd}}_{i,t+h})\}$; per-side cost $c_h$; round-trip flag $\rho \in \{0,1\}$; annualisation $a_h$
\Ensure per-rebalance P\&L frame + annualised summary stats
\For{each rebalance date $t$}
  \State rank $U_t$ by $s_{i,t}$, ties broken stably
  \State assign terciles $q_{i,t} \in \{\text{lo}, \text{mi}, \text{hi}\}$ via Eq.~\eqref{eq:weights}
  \State compute $w^{\MOM}_{i,t}$ and $w^{\REV}_{i,t} = -w^{\MOM}_{i,t}$
  \State gross $\Pi^{\MOM}_t \gets \sum_i w^{\MOM}_{i,t}\, r^{\text{fwd}}_{i,t+h}$
  \State gross $\Pi^{\REV}_t \gets -\Pi^{\MOM}_t$ \Comment{exact mirror}
\EndFor
\State pivot $w^{\MOM}$ to wide (date $\times$ ticker) frame $W$
\If{$\rho = 1$} \Comment{intraday round-trip: H1, H2}
  \State $\tau_t \gets 2$ for every rebalance
\Else \Comment{carry-through turnover: H3--H6}
  \State $\tau_t \gets \sum_i \bigl|W_{t,i} - W_{t{-}h,i}\bigr|$; set $\tau_0 = \sum_i |W_{0,i}|$
\EndIf
\State cost drag $\kappa_t \gets c_h \cdot \tau_t$
\State net $\Pi^{\MOM,\text{net}}_t \gets \Pi^{\MOM}_t - \kappa_t$; same for REV
\State \textbf{invariant:} $\Pi^{\MOM}_t + \Pi^{\REV}_t = 0$ (mirror)
\State \textbf{invariant:} $\tau_t \in [0, 2]$
\State annualise: $\hat\mu_{\text{ann}} \gets \bar{\Pi^{\text{net}}}\,a_h$, $\hat\sigma_{\text{ann}} \gets \mathrm{sd}(\Pi^{\text{net}})\sqrt{a_h}$
\State Sharpe $\gets \hat\mu_{\text{ann}}/\hat\sigma_{\text{ann}}$
\State $t \gets \bar{\Pi^{\text{net}}} / \mathrm{sd}(\Pi^{\text{net}})\cdot \sqrt{N}$
\If{H6} $t \gets $ Newey--West-adjusted $t$ via Eq.~\eqref{eq:nw}
\EndIf
\State \Return \{gross, net, $\tau$\} per date; annualised summary
\end{algorithmic}
\end{algorithm}

The mirror invariant ($\Pi^{\MOM} + \Pi^{\REV} = 0$ to machine precision)
is a self-check: any non-zero sum implies a bug in the weight assignment.
We assert it after every run.

%======================================================================
\section{The Six Horizons}
\label{sec:six}
%======================================================================

\begin{table*}[t]
\centering
\caption{Six horizons, three data sources. $a_h$ is the annualisation
factor (rebalances per year); $c_h$ is the per-side cost. Rebalance
``bar'' means every 30 minutes of the regular session; ``day''
means end-of-day; ``month-end'' means the last trading day of each
calendar month.}
\label{tab:horizons}
\small
\begin{tabular}{@{}l l l r r l l@{}}
\toprule
\textbf{$h$} & \textbf{Horizon} & \textbf{Dataset} & \textbf{$a_h$} &
\textbf{$c_h$ (bps)} & \textbf{Rebalance} & \textbf{Classical prior} \\
\midrule
H1 & 30 min          & TAQ       & 2{,}772 & 3.0 & every bar          & Microstructure reversal / momentum~\cite{roll1984,gao2018} \\
H2 & Intraday AM/PM  & iid\_ms   & 252     & 1.5 & open + close       & Intraday momentum~\cite{gao2018} \\
H3 & 1 day           & CRSP      & 252     & 1.5 & daily close        & One-day reversal~\cite{lehmann1990} \\
H4 & 5 day (weekly)  & CRSP      & 52      & 1.5 & ISO week-end       & Weekly contrarian~\cite{lo1990} \\
H5 & 21 day          & CRSP      & 12      & 1.5 & month-end          & One-month reversal~\cite{jegadeesh1990} \\
H6 & 126 day (6 mo)  & CRSP      & 12      & 1.5 & month-end (overlap)& JT momentum~\cite{jt1993} \\
\bottomrule
\end{tabular}
\end{table*}

\subsection{Rebalance-date convention}

H3 rebalances every trading day; the signal is today's total return
and the forward is tomorrow's. H4 uses the last trading day of each
ISO calendar week (a holiday-adjusted Friday proxy); H5 uses the last
trading day of each calendar month; H6 uses month-ends with a 126-day
forward (so monthly returns overlap by five months and need
Newey--West, \S\ref{sec:nw}). H2 has one rebalance per day with entry
at $\sim$10:00 ET (\texttt{mid\_after\_open}) and exit at 16:00
(\texttt{CPrc}); we drop the 19 NYSE early-close days
(Black Friday, Christmas Eve, July 3) to avoid compressed-session
contamination of the rest-of-day window.

\paragraph{Why these six horizons.} The six span roughly nine
orders of magnitude in time ($\sim\!10^{3}$~s at H1 to $\sim\!10^{7}$~s
at H6). Each corresponds to a different mechanism named in the
literature: Roll's bid-ask bounce at H1, overnight-news
absorption at H2, Lehmann's daily liquidity provision at H3,
Lo--MacKinlay's weekly lead-lag at H4, Jegadeesh's
one-month reversal at H5, and Jegadeesh--Titman's intermediate
momentum at H6. Picking fewer horizons would leave gaps that
allow a reader to ask ``what about \emph{x}\,?'' for some
intermediate $x$; picking more would add multiple-testing cost
without adding mechanism. Six is the smallest set that covers
the dominant mechanisms cited in the literature and gives a
contiguous horizon axis to plot against.

\paragraph{Why the Jegadeesh--Titman skip-1-month convention on
H6.} Without the skip, the H6 signal window $[t-125, t]$ and the
H5 signal window $[t-20, t]$ would share 21 of the 126 days of
lookback. Any success at H6 would then be partially attributable
to whatever signal H5 carries on its last 21 days, and the two
horizons could not be interpreted as independent. The skip
$[t-125, t-21]$ isolates the intermediate-momentum channel from
the short-term reversal channel cleanly.

\begin{figure}[ht]
\centering
\includegraphics[width=\linewidth]{icm/figures/windowing_diagram.pdf}
\caption{Signal (green) and forward-hold (red/orange) windows per
horizon, in trading days relative to rebalance date $t$. H6 is the
only overlapping horizon: its 126-day forward extends past the next
monthly rebalance, and its signal uses a 5-month window ending 21
days \emph{before} $t$ (Jegadeesh--Titman skip-1-month, shown
hatched). H1 and H2 windows are sub-daily and are shown with
annotations rather than at physical scale.}
\label{fig:windowing}
\end{figure}

Figure~\ref{fig:windowing} makes the alignment unambiguous. The
key diagnostic is that for every horizon, $W_{\text{past}}$ ends at
or before $t$ and $W_{\text{fwd}}$ starts at or after $t$. No window
contains $t$ in its interior; no window contains a point from the
other side. Look-ahead leakage would require breaking this
separation, and the look-ahead audit in \S\ref{sec:robust} confirms
it is not broken: shifting the signal back by one rebalance kills
the H2 REV result entirely (Sharpe $+2.53 \to -1.52$).

\subsection{Panel integrity and signal dispersion}
\label{sec:integrity}

\begin{table*}[t]
\centering
\caption{Per-horizon panel integrity and cross-sectional signal
moments. ``Modal $N$'' is the most common number of stocks present
per rebalance (target 30); ``short-days'' counts rebalance dates
with $<$30 names (WRDS data-quality gaps, ticker transitions, and
--- for H1 --- the TAQ pull that is 39 tickers not 40). ``Turn.'' is
the mean per-rebalance turnover. $\bar{s}$ and $\sigma_s$ are the
pooled mean and standard deviation of the signal distribution (bps);
XS-$\sigma_s$ is the mean per-date cross-sectional standard
deviation --- the actual dispersion the tercile sort ranks against.
Kurtosis is \emph{excess} over Gaussian.}
\label{tab:integrity}
\scriptsize
\input{icm/tables/signal_moments}
\end{table*}

Table~\ref{tab:integrity} is the pre-results panel audit. Three
things to notice.

\emph{Modal names per rebalance is 30 on every CRSP and iid horizon},
with a small number of short days (3--21) attributable to known
ticker-chain boundaries or vendor-side $\texttt{mid\_after\_open}$
gaps. H1 is the exception at modal~29 because the TAQ pull predates
the supplementary Goldman Sachs pull that fixed CRSP and iid; H1
therefore operates on a 39-ticker universe (covered in \S\ref{sec:h1-note}).

\emph{Cross-sectional signal dispersion grows roughly with
$\sqrt{h}$} across CRSP horizons: 122~bps at H3 (1-day),
274~bps at H4 (5-day), 565~bps at H5 (21-day), 1{,}321~bps at H6
(126-day skip-21). The intraday horizons H1 and H2 cluster at the
bottom, as expected (only $\sim$12--22 bps of cross-sectional
spread per bar / first-30). A tercile long--short pays in
proportion to this dispersion: more dispersion means more room
between the top and bottom terciles' mean signals.

\emph{Every signal has heavy tails}: excess kurtosis is positive at
every horizon, ranging from 2.35 at the monthly/semi-annual
horizons to 32.7 at H2 (intraday first-30 moves are dominated by
news-day fat tails). This motivates the rank-based portfolio
construction: a mean--variance or inverse-vol weighting would load
heavily on a handful of outlier days, while the rank sort assigns
every name to a tercile regardless of how extreme its signal is.

\paragraph{H1 universe note.}
\label{sec:h1-note}
The TAQ pull on MIT Engaging includes 39 of 40 DJ30 ever-members
(Goldman Sachs missing). H1 therefore runs on a 29-name universe
every day. With terciles of size $(9, 11, 9)$ or $(10, 9, 10)$ the
portfolio is still dollar-neutral and gross-one, but the
$\sqrt{9/10} = 0.95\times$ noise penalty on the smaller legs
modestly widens per-date P\&L dispersion. This is a cosmetic concern:
H1 is dominated by transaction costs by roughly three orders of
magnitude (break-even 0.004 bps per side), so the missing ticker
changes none of the qualitative result.

\subsection{H1 bar aggregation on TAQ}

The TAQ raw files are $\sim$93~GB gzipped. A SLURM array of ten
per-year tasks runs on the Engaging cluster. Each task streams a
year's gzip in 5M-row chunks, applies trade filters (regular session,
\texttt{TR\_CORR=00}, positive price and size, \texttt{SYM\_ROOT}
in the ever-member set, empty \texttt{SYM\_SUFFIX} to drop
preferreds, and a \texttt{TR\_SCOND} whitelist that excludes
O/6/L/Z/T/W/U/I modifiers~\cite{nysetaq}), assigns each surviving
trade to one of 13 half-hour bars ($k \in\{1,\ldots,13\}$, $k{=}1$
for 09:30--10:00), drops $k{=}1$ and $k{=}13$ to avoid auction
effects, and aggregates to OHLCV+VWAP bars. The pipeline processes
10.3 billion raw trades into 1.01 million surviving bars in
$\sim$2~hours wall-clock. On early-close days the grid collapses
to 7 bars of which we keep $k{=}2$--$6$. A dedicated pytest suite
(35 cases) covers the filter, the bar-index math, and the
end-to-end aggregation on a tiny synthetic gzip.

Bar returns are log-returns; signal at bar $k$ is
$r_{k{-}1}$ (previous bar) and forward is $r_{k{+}1}$. All shifts
are strictly within (ticker, trading\_date) so we never cross the
overnight boundary. Of 92{,}535 (ticker, date) cells audited,
92{,}534 have exactly the expected bar count (11 regular / 5
early-close), and daily VWAPs reconstructed from bars match
CRSP's \texttt{DlyPrcVol}/\texttt{DlyVol} within 0.3\% on 5 random
sample days.

\paragraph{Why drop bars $k\!=\!1$ and $k\!=\!13$.} The opening
auction concentrates the overnight order imbalance into a single
print at 09:30; the first continuous 30 minutes inherit that print
and any residual imbalance still being worked. Prices during that
window are dominated by imbalance-absorption dynamics, not by
normal continuous trading. Symmetrically, the last 30 minutes
before 16:00 are dominated by MOC (market-on-close) orders
routing into the closing auction. Both auction-adjacent windows
are distinct price-formation regimes from the 10:00--15:30 bulk
session we want to study, so we drop them at the aggregation
stage rather than try to correct for them post hoc. The loss is
one bar of 11 at each end --- $\sim$18\% of the continuous
session, but the 18\% that carries the most non-signal noise.

%======================================================================
\section{Headline Comparison}
\label{sec:headline}
%======================================================================

\begin{figure*}[t]
\centering
\includegraphics[width=\linewidth]{icm/figures/figure3_horizons.pdf}
\caption{(a) Annualised net Sharpe by horizon and direction.
(b) $t$-statistics with the two-sided Bonferroni threshold
$t^{*}\approx 2.87$ (dashed). H1 is off-scale (cost-dominated) and
is annotated textually. Only H2 REV clears Bonferroni; every other
row is $\text{N.S.}$ at the family-wise level.}
\label{fig:horizon}
\end{figure*}

Figure~\ref{fig:horizon} is the central exhibit.
Table~\ref{tab:detail} gives the full per-horizon numbers. Four
passes through the table:

\emph{H1 (30-min).} Gross mean per bar is $6\times10^{-7}$
(essentially zero). Cost drag at $3\text{~bps/side}\times1.35~\text{turnover}\times 2{,}772$ is
$\sim 112\%$ annualised. Both directions converge to net Sharpe
$\approx -34$ and MDD $\approx -100\%$. Break-even per-side cost
is 0.004~bps, three orders of magnitude below plausible
execution cost. The H1 regime is one where transaction costs
dominate any cross-sectional signal at this frequency --- a useful
null and a reminder that rebalancing cadence matters as much as
signal direction.

\emph{H2 (intraday).} \textbf{REV earns net Sharpe $+2.53$ at
$t = +7.86$,} with max drawdown only $-4.3\%$. This is the first row
of the table to clear Bonferroni and the only one that does so at
cost assumptions we consider reasonable (\S\ref{sec:cost-sens}). Its
economic anchor is a per-ticker regression of rest-of-day return
on the first-30 signal: cross-sectional mean $\beta = -0.64$
$(t_{\text{cross-sec}} = -5.27)$, 28 of 40 tickers individually
significantly negative, 0 of 40 significantly positive. The sign
differs from Gao et al.'s~(2018) index-level finding on S\&P~500
futures (positive mean $\beta$). We do not view this as a
contradiction: the two tests have different designs (individual
large-cap names vs.\ the S\&P~500 index, cross-sectional ranking
vs.\ time-series return autocorrelation, 2016--2025 vs.\ an
earlier window). On individual DJ30 names in our sample, the
first-30 move appears to overshoot and revert rather than
persist.

\emph{H3--H5 (daily--monthly).} All three fall inside the noise
band. H3 REV has net Sharpe $-0.74$ (cost-dominated); H4 REV is
$+0.04$ (within noise); H5 REV is $-0.22$ (also within noise, and
opposite in sign to Jegadeesh's~(1990) broad-universe estimate,
which is not surprising on a 30-name large-cap subset). Nagel~\cite{nagel2012}
documented that short-term reversal weakened substantially between
the 1990s and the 2000s; our DJ30 CRSP-horizon results are
consistent with that trend extending into the 2016--2025 window.
More broadly, the Adaptive Markets Hypothesis~\cite{lo2004} offers
a natural framing: the strength of a given anomaly varies with
the population of market participants and their technology, so the
absence of a 1990s-vintage effect in a 2016--2025 sample is
unsurprising.

\emph{H6 (6-month).} The only CRSP row with directionally classical
behaviour: MOM net Sharpe $+0.82$ with naive $t = +2.43$. But the
naive $t$ is inflated by overlap; NW $q{=}6$ brings it to
\textbf{$+1.37$}, below even uncorrected 1.96 two-sided, let alone
Bonferroni. H6 MOM is the best-available momentum candidate;
because no MOM horizon passes the family-wise test we present it
with the caveat that its significance is not confirmed.

%======================================================================
\section{H2 REV: The Winning Strategy}
\label{sec:h2}
%======================================================================

\subsection{Economic mechanism}

H2 REV's signal is the return from the opening print to the
$\sim$5-minute VWAP midprice just after the open. Its forward is the
return from that midprice to the official close. The portfolio is
long the bottom tercile (``openers that fell'') and short the top
tercile (``openers that rose''), equal-weighted, dollar-neutral,
unwound every close.

Why does this work? Two mechanisms jointly match the data:

\paragraph{Overnight-news overshoot.} Overnight news is absorbed at
the NYSE/NASDAQ opening auction and in the first minutes of
continuous trading. For DJ30 names (which have tight quoted spreads
and deep liquidity) the open-to-first-30 move tends to \emph{overshoot}
the fundamental price adjustment as market-makers absorb one-sided
flow. The subsequent six hours then mean-revert that overshoot as
liquidity providers work their inventory back to target.

The microstructure rationale is an inventory cycle. Any order at
the open is demanding immediate liquidity --- by construction, there
is no pre-09:30 book to rest against. Market-makers accept that
flow but widen quotes to compensate for the inventory risk of
holding a one-sided position. The first 30 minutes' price therefore
reflects not just the fundamental news but also an inventory
premium paid \emph{to} the market-maker. Over the rest of the day,
as the dealer works inventory off through natural two-sided flow,
the premium decays. That decay appears as a mean-reversion in the
observed mid-price from $\sim$10:00 to 16:00 --- the pattern that
H2 REV is designed to capture.

\paragraph{Cross-sectional ranking cancels the market move.} The
long--short is dollar-neutral and long-10/short-10; its P\&L depends
on \emph{relative} first-30 moves across the 30 names that morning,
not on the absolute market direction. This is why the classical
Gao intraday-momentum story (which is a time-series, per-instrument
test on S\&P~500 futures) can be positive while the cross-sectional
rank strategy on DJ30 individual names is negative: the two tests
are asking different questions.

\begin{figure}[ht]
\centering
\includegraphics[width=\linewidth]{icm/figures/equity_h2_rev.pdf}
\caption{H2 REV cumulative wealth (gross and net of $1.5~\text{bps/side}$)
and drawdown panel. 10 years of daily rebalances, $N = 2{,}431$.}
\label{fig:equity-h2}
\end{figure}

\subsection{Equity curve and drawdown}

Figure~\ref{fig:equity-h2} shows the wealth curve and drawdowns.
Gross wealth grows to $5.0\times$ over 10 years; net of costs to
$2.6\times$. Drawdowns are remarkably shallow for a leverage-one
long--short: the deepest is $-4.3\%$, with no month ever posting a
loss greater than 3\%. The intraday-only mechanic --- every day
starts flat --- eliminates overnight-gap tail risk that would
otherwise dominate a strategy with tercile-based ranks.

\subsection{Factor regression}

Table~\ref{tab:factors} reports OLS regressions of daily net P\&L
on the Fama--French 5 factors~\cite{ff2015} plus the
Carhart-style momentum factor~\cite{carhart1997}. Every factor
$t$-statistic on H2 REV is below 1.86 in absolute value. The
intercept is $\hat\alpha = +0.041\%$ per day $= +10.3\%$
annualised, with $t = +7.82$, and $R^{2} = 0.003$. The return
stream is largely unexplained by standard daily risk factors.
This is what we would expect for an intraday strategy: positions
are flat by 4:00~PM each day, so daily factor movements cannot load
onto the H2 P\&L through overnight exposure.

\begin{table}[ht]
\centering
\caption{Fama--French 5 + Mom factor regression on daily net P\&L of
H2 REV (OLS) and monthly net P\&L of H6 MOM (NW $q{=}6$). $\alpha$
is annualised; factor coefficients are decimal. H2 REV shows no
significant loading on any of the six factors, consistent with its
intraday-only structure.}
\label{tab:factors}
\small
\input{icm/tables/factor_regression}
\end{table}

\subsection{Per-ticker contribution}

A useful sanity check for any long--short on 30 names is whether the
return is concentrated in a handful of tickers. We decompose the
annualised REV return by summing per-ticker per-day contribution
$w^{\REV}_{i,t}\cdot r^{\text{fwd}}_{i,t}$.

\begin{figure}[ht]
\centering
\includegraphics[width=\linewidth]{icm/figures/h2_ticker_contribution.pdf}
\caption{Per-ticker contribution to H2 REV's annualised return.
Top-5 contributors produce 40\% of the total; the effect is spread
across the cross-section.}
\label{fig:tickers}
\end{figure}

Figure~\ref{fig:tickers} shows the top-5 tickers contribute 40\% of
the total annualised return, with the remaining 60\% spread across
the other 35 names. The contribution is reasonably diffuse for a
30-name universe, which argues against a single-name driver
interpretation (e.g.\ ``the strategy is just long AAPL after
dips'').

%======================================================================
\section{H6 MOM: Fallback Momentum}
\label{sec:h6}
%======================================================================

H6 MOM ranks stocks by their 105-trading-day return ending 21 days
ago (the skip-1-month JT convention) and holds the long--short
tercile portfolio for 126 trading days, rebalancing monthly. The
signal and forward window overlap across consecutive rebalances by
$105/126 = 83\%$ of length, hence the Newey--West correction.

\begin{figure}[ht]
\centering
\includegraphics[width=\linewidth]{icm/figures/equity_h6_mom.pdf}
\caption{H6 MOM cumulative wealth and drawdown; 105 monthly
rebalances. Net Sharpe $+0.82$ (naive $t = +2.43$; NW $q{=}6$
$t = +1.37$, below Bonferroni).}
\label{fig:equity-h6}
\end{figure}

Figure~\ref{fig:equity-h6} shows H6 MOM's equity curve. The strategy
grows net wealth to $3.1\times$ but pays for it with a $-48\%$
drawdown in 2020. Moskowitz--Grinblatt~(1999)~\cite{mg1999} argue
that most of cross-sectional momentum at the stock level is
industry-level momentum; on a 30-name universe with one or two
stocks per industry, the industry channel is attenuated, and the
effect we observe (positive but below significance) is consistent.

The FF5+Mom regression (Table~\ref{tab:factors}, right column)
confirms H6 MOM is only weakly explained by the standard daily
momentum factor (loading $+0.04$, $t = +0.32$) --- the monthly
Carhart momentum is ``industry-plus-large'', while our DJ30
construction is ``top-10-of-30 minus bottom-10.'' The residual
alpha $t = +1.64$ (NW) is below significance at every standard
threshold. We present H6 MOM as the ICM's best-available momentum
candidate while clearly noting the failure of the Bonferroni test.

%======================================================================
\section{Statistical Properties of H2 REV}
\label{sec:stats}
%======================================================================

\begin{figure}[ht]
\centering
\includegraphics[width=\linewidth]{icm/figures/h2_return_distribution.pdf}
\caption{(a) Daily P\&L distribution in bps with fitted normal,
VaR$_{95}$ and ES$_{95}$ annotations. (b) QQ plot vs.\ normal.
Skew is near zero; excess kurtosis is $+3.4$ --- heavier tails
than Gaussian but symmetric.}
\label{fig:dist}
\end{figure}

Figure~\ref{fig:dist} characterises the daily return
distribution. $\mu = +4.1$~bps, $\sigma = 25.6$~bps; skewness
$-0.02$ (essentially symmetric); excess kurtosis $+3.4$ (fat tails
compared to Gaussian, as expected for any daily return series).
The 5\% value-at-risk is $-37$~bps; expected shortfall conditional
on the worst 5\% is $-54$~bps. The QQ plot shows the tails bow
outward symmetrically --- the tails are heavier than normal in
both directions, without a bias toward the downside.

The symmetry is economically informative. Merger-arbitrage and
short-volatility strategies typically carry \emph{left}-skew: most
days are small gains, a small number of days are large negative
outliers, and the Sharpe ratio tends to understate tail risk. H2
REV has no such asymmetry --- the worst 5\% of days
($-54$~bps expected shortfall) roughly mirrors the best 5\%. That
is more consistent with the market-maker inventory story above
than with a ``selling insurance'' interpretation of the edge,
though we would not claim to have ruled the latter out
definitively from a 10-year sample.

\begin{figure}[ht]
\centering
\includegraphics[width=\linewidth]{icm/figures/h2_rolling_sharpe.pdf}
\caption{Rolling 252-day Sharpe ratio for H2 REV. The window-average
never dips below $+0.3$ after 2017; full-sample Sharpe is the dashed
line.}
\label{fig:rolling}
\end{figure}

Figure~\ref{fig:rolling} plots a rolling one-year Sharpe. The
strategy's Sharpe has ranged between $+0.4$ and $+5.5$ across
overlapping 252-day windows, with the upper end concentrated in
2022--2025. The rolling window never goes negative after 2017.

\begin{figure}[ht]
\centering
\includegraphics[width=\linewidth]{icm/figures/h2_autocorr.pdf}
\caption{Autocorrelation of daily H2 REV net P\&L out to lag 30,
with 95\% confidence band at $\pm 1.96/\sqrt{N}$.}
\label{fig:acf}
\end{figure}

Figure~\ref{fig:acf} shows the autocorrelation of H2 REV's daily
P\&L. Three lags are marginally outside the 95\% band; 27 are
inside. The P\&L series is close to i.i.d., consistent with an
intraday strategy where the information set re-shuffles each
morning. The absence of meaningful autocorrelation justifies
treating the Sharpe annualisation factor of 252 as correct (no
overlap correction required). Should a more conservative
inference be desired, the stationary bootstrap of
Politis--Romano~\cite{pr1994} with average block length 20 and
$B = 2{,}000$ replicates produces a two-sided 95\% confidence
interval for the annualised Sharpe of $[1.82, 3.27]$, comfortably
above zero.

%======================================================================
\section{Robustness}
\label{sec:robust}
%======================================================================

The robustness suite interrogates three prior objections an
Investment Committee will raise against H2 REV: is it eaten by
costs, is it a microstructure artefact, is it a regime fluke?
Each exhibit below attacks a \emph{specific} null that a skeptical
committee would name; each ends with a one-sentence verdict so the
reader can check whether the null was in fact defeated.

\subsection{Cost sensitivity and break-even}
\label{sec:cost-sens}

\begin{figure}[ht]
\centering
\includegraphics[width=\linewidth]{icm/figures/h2_cost_sensitivity.pdf}
\caption{(a) H2 REV net Sharpe as per-side cost rises. (b) $t$-stat
vs.\ cost. Break-evens: Sharpe$=0$ at $3.54~\text{bps/side}$;
Bonferroni at $2.80~\text{bps/side}$.}
\label{fig:cost}
\end{figure}

The engine charges 1.5 bps per side $\times$ 2.0 turnover = 3~bps
round-trip per day $\approx 7.5\%$ annualised cost drag.
Figure~\ref{fig:cost} sweeps the per-side cost from 0 to 20~bps.
Two break-evens matter: at $c = 3.54$~bps/side the net Sharpe
crosses zero; at $c = 2.80$~bps/side the NW-corrected $t$ drops
below $t^{*}$. Realistic execution on DJ30 is at half-spread
plus impact $\sim$1.0--1.5~bps/side, leaving a factor-of-two
margin before Bonferroni is breached.

\subsection{Roll bid-ask-bounce diagnostic}

\begin{figure}[ht]
\centering
\includegraphics[width=\linewidth]{icm/figures/h2_beta_vs_spread.pdf}
\caption{Per-ticker Gao $\beta$ vs.\ median quoted spread. If H2 REV
were a Roll~(1984) artefact, $\beta$ would scale with spread; the
observed $\rho = -0.10$ ($R^{2} = 0.01$) rules this out.}
\label{fig:beta-spread}
\end{figure}

Roll~(1984)~\cite{roll1984} shows that an observed price series
with noisy bid-ask bounce has a mechanical first-lag
autocorrelation of $\rho_1 = -(s/2\sigma)^2$, where $s$ is the
effective spread and $\sigma$ is the fundamental return standard
deviation. For typical DJ30 values ($s \approx 1$--$8$~bps,
$\sigma \approx 100$--$200$~bps per day), this predicts a $\beta$
contribution of $10^{-4}$ to $10^{-3}$ --- six to seven orders of
magnitude below the $-0.64$ we observe. Figure~\ref{fig:beta-spread}
confirms empirically: the cross-sectional correlation between
per-ticker $\beta$ and median spread is $-0.10$ ($t = -0.62$,
$R^{2} = 0.010$). If the mechanism were pure spread-noise reversal,
$\rho$ would approach $-1$; at $-0.10$ the bid-ask bounce
explains 1\% of cross-sectional $\beta$ variation. Restricting the
backtest to the tightest-spread half of the universe (20 tickers,
average spread 1.79~bps) still yields net Sharpe $+1.23$ at
$t = +3.81$, clearing Bonferroni on that subset alone.

\subsection{Time and regime stability}

\begin{figure}[ht]
\centering
\includegraphics[width=\linewidth]{icm/figures/h2_by_year.pdf}\\[2pt]
\includegraphics[width=\linewidth]{icm/figures/h2_by_vix_quintile.pdf}
\caption{(top) H2 REV Sharpe and annualised return by calendar year.
(bottom) Sharpe and within-bucket $t$ by VIX-prior-close quintile.
10/10 calendar years are positive; 5/5 VIX regimes are positive.}
\label{fig:regime}
\end{figure}

Figure~\ref{fig:regime} breaks the daily P\&L by calendar year
(top) and by prior-day VIX quintile (bottom). The strategy has a
positive Sharpe in \emph{every} year from 2016 to 2025 (mid-year)
and in every VIX regime from the calmest (quintile~1, median
VIX 12) to the most turbulent (quintile~5, median VIX 28).
The effect has strengthened post-2020 (Sharpes of 2.2--5.7 across
2020--2025 vs.\ 0.4--1.3 across 2016--2019), consistent with the
hypothesis that intraday reversal scales with cross-sectional
first-30 dispersion, which has been higher post-2020.

\subsection{Survivorship, look-ahead, and variance ratio}

\paragraph{Survivorship.} Restricting H2 to the 24
always-present DJ30 tickers (full 2{,}514-day tenure) drops Sharpe
from 2.53 to 2.37 and $t$ from 7.86 to 7.35 --- still
Bonferroni-significant on the restricted subset. The strategy is
not driven by index-rotation winners.

\paragraph{Look-ahead audit.} Shifting the H2 signal back by one
rebalance (using yesterday's first-30 to predict today's
rest-of-day) changes Sharpe from $+2.53$ to $-1.52$. The signal
is therefore correctly ex-ante: only same-day information
carries predictive power. A complete loss of Sharpe in a
look-ahead test would be a red flag for signal leakage; the
sign-flip we observe is consistent with an honest,
contemporaneous-information strategy.

\begin{figure}[ht]
\centering
\includegraphics[width=\linewidth]{icm/figures/vr_term_structure.pdf}
\caption{Lo--MacKinlay variance-ratio term structure for 40 DJ30
ever-members. $\mathrm{VR}(q) < 1$ is the reversal signature.
Median drops from $0.95$ at $q{=}2$d to $0.56$ at $q{=}252$d; the
dominant regime is reversal across every horizon tested.}
\label{fig:vr}
\end{figure}

\paragraph{Variance ratio term structure.} Figure~\ref{fig:vr}
computes $\VR(q)$ per ticker for $q \in \{2,5,21,63,126,252\}$d
using the Lo--MacKinlay~(1988)~\cite{lm1988}
heteroskedasticity-robust framework --- the standard test in this
literature. The median is monotonically below 1 at every horizon:
$0.95$ at $q{=}2$, $0.88$ at $q{=}21$, $0.63$ at $q{=}126$,
$0.56$ at $q{=}252$. $92.5\%$ of ticker-years have
$\VR(2) < 1$; $95.0\%$ have $\VR(126) < 1$. The time-series
evidence for per-stock reversal across horizons is strong and
aligns with the classical Lo--MacKinlay finding on broader
universes.

\paragraph{How the time-series and cross-sectional pictures fit
together.} $\VR(q)$ is a per-stock measure: it summarises
$\sum_{k=1}^{q-1}(1-k/q)\rho_k$, the weighted sum of own-lag
autocovariances of a single return series. Our tercile
long--short asks a distinct question: given the cross-section of
today's signals, are the names at the extremes relatively more
or less attractive going forward? The two questions are related
but not equivalent. For example, a common market factor can
contribute positively to $\rho_k$ of every stock without creating
any cross-sectional ranking effect (the factor cancels in a
dollar-neutral long--short). Conversely, cross-sectional ranking
effects can exist without per-stock autocorrelation if the
ranking is driven by slowly-varying characteristics. On DJ30 we
see both phenomena: strong per-stock reversal in the VR measure
\emph{and} a profitable cross-sectional strategy at H2, but not at
most CRSP horizons. The Lo--MacKinlay VR evidence is therefore
complementary to our headline, not competing with it: both views
are consistent with a reversal-rich regime, and both help
characterise the environment in which H2 REV operates.

%======================================================================
\section{Headline Table}
\label{sec:table}
%======================================================================

\begin{table*}[t]
\centering
\caption{Per-horizon detailed results. $c$ is the assumed per-side
cost; turn.\ is the average turnover per rebalance; $c^{*}$ is the
break-even per-side cost (gross mean divided by average turnover).
${}^{*}$ indicates rows where $t$ clears the Bonferroni threshold
$t^{*}\approx 2.87$ on the correct (positive) side.}
\label{tab:detail}
\small
\input{icm/tables/per_horizon_detail}
\end{table*}

Table~\ref{tab:detail} is the comprehensive per-horizon table.
Only H2 REV is starred. The break-even column is a useful
alternative lens: H6 MOM's $c^{*} = 179$~bps per side suggests
the strategy is cost-robust if one accepts the gross signal, but
of course the gross signal fails Bonferroni. H5 MOM's $c^{*} =
6.5$~bps per side is the second most cost-robust row, but its
naive Sharpe is $+0.14$ with $t = +0.42$ --- not actionable.

%======================================================================
\section{Implementation and Risk}
\label{sec:impl}
%======================================================================

\subsection{Capacity}

H2 REV holds each tercile for six hours and trades at
$\sim$10:00 ET and 16:00 ET, when DJ30 liquidity is deepest.
Average daily dollar volume across the 30 names in 2024--2025 is
$\sim \$3$B per ticker; a $1\%$ share-of-volume participation
cap gives $\sim \$30$M executable per name per 30-minute
window. With 10 names per leg and the portfolio dollar-neutral,
gross book capacity at $1\%$ participation is $\$300$M one-sided
or $\$600$M total. Higher participation is feasible for DJ30 but
will move the average execution cost upward from the $\sim$1.5~bps
quoted-spread anchor toward the $\sim$2--3~bps range, which the
cost-sensitivity curve (\S\ref{sec:cost-sens}) shows the strategy
still tolerates.

\subsection{Kill switches}

Three rules protect against live-trading regime changes:

\begin{enumerate}[leftmargin=*]
  \item \emph{30-day rolling Sharpe $< 0$ for 20 consecutive
        trading days.} The rolling Sharpe has never been negative
        for more than 8 days in the backtest; triggering this
        rule would represent a statistically novel regime.
  \item \emph{Drawdown $> -5\%$.} Historical max is $-4.3\%$.
  \item \emph{First-30 dispersion falls below the $5$-th percentile
        of 2016--2019 levels.} The strategy mechanically pays in
        proportion to dispersion; if dispersion collapses,
        expected alpha collapses too.
\end{enumerate}

%======================================================================
\section{Limitations}
\label{sec:limits}
%======================================================================

\emph{Small cross-section.} DJ30 has 30 names per day. Tercile
sizes of 10 are small; per-date rank noise is higher than in a
500-name universe. The \emph{time}-series length (2{,}431
trading days) compensates, giving $t{=}7.86$, but the effect at
any single date is noisy.

\emph{Calendar trim.} CRSP ends 2025-09-30; iid and TAQ extend
to 2025-12-31. We trim all series to 2025-09-30 for a shared
calendar. Q4 2025 data is available for H2 and would add a few
more months to the sample; we exclude it to preserve the
horizon-comparison consistency.

\emph{Cost model conservatism.} The $1.5~\text{bps}$/side
assumption is conservative for DJ30 large caps, which usually
trade at half-spread of 0.5--1.0~bps. A more aggressive cost
of 1.0~bps/side would raise H2 REV's net Sharpe to $+3.15$
(Figure~\ref{fig:cost}); we nonetheless report the conservative
number in all headline figures.

\emph{Overnight reversal not tested.} The strategy closes every
day at 16:00. An overnight extension (hold positions through
the close to the next open) could add or subtract alpha; we do
not test it here.

\emph{Single-vintage backtest.} 2016--2025 covers two major
regimes (post-QE rise 2016--2019, post-COVID high-vol era
2020--2025). An earlier vintage (e.g.\ 2004--2015) would be a
better robustness check for regime-persistence; daily
iid\_ms coverage does not go back that far.

%======================================================================
\section{Conclusions}
\label{sec:conclusions}
%======================================================================

A horizon-matched comparison of tercile long--short strategies on
the DJ30 from 2016 to 2025 produces a single Bonferroni-significant
strategy out of 12 tests: an intraday mean-reversion strategy that
buys the morning's worst-opening DJ30 names and sells the
best-opening ones, held to the close. Net of $1.5$~bps/side
costs, the strategy earns $+10.3\%$ annualised at a Sharpe of
$+2.53$, with a maximum drawdown of $-4.3\%$ and a $t$-statistic
of $+7.86$, comfortably above the Bonferroni threshold of
$t^{*}\!\approx 2.87$ for 12 tests. A FF5+Mom factor regression
delivers no significant factor loadings and a $+10.3\%$ annualised
alpha ($t = +7.82$), consistent with the strategy's intraday-only
structure.

Every other horizon falls short of the Bonferroni threshold for
reasons we can identify: transaction costs dominate H1 (bar-level
rebalance) and H3 (daily rebalance); H5 and H4 produce results
within noise on DJ30 large caps; H6 Jegadeesh--Titman momentum
has the right sign but its overlap-adjusted $t$-stat
(NW $q{=}6$, $t = +1.37$) does not clear family-wise significance.
The IC's recommendation is \textbf{REV-H2} as the primary
mean-reversion strategy; H6 MOM is presented as the
best-available candidate for the momentum family, with the caveat
that its significance is not confirmed at the family-wise level.

Taken together, the sweep offers a single-universe,
single-methodology reading of which classical horizons generate
tradable edge on DJ30 in 2016--2025 and which do not --- a reading
we view as a small complement to the broader-universe evidence in
the literature, not a replacement for it.

%======================================================================
\section*{LLM Usage and Attribution}
%======================================================================

Code scaffolding, debugging, data-quality diagnosis, and document
formatting were done with Claude (Anthropic) under human 
supervision across a series of numbered
attribution logs kept alongside the code. Computations were performed on MIT
Engaging (TAQ aggregation, $\sim$2~hours wall on a 10-task SLURM
array) and on a laptop (everything else, $\sim$30~seconds to
regenerate every figure and table from cached P\&L series).

%======================================================================
\appendix
%======================================================================

\section{Data-quality ledger}
\label{app:ledger}

Table~\ref{tab:ledger} enumerates every integrity check the pipeline
asserts before any backtest is run. Green rows pass on every build;
red rows would halt the pipeline. The H1-specific row about 39 vs.\
40 tickers is the only amber: it is a known data gap documented
inline rather than a check failure.

\begin{table}[ht]
\centering
\caption{Pre-backtest integrity ledger. Every row runs as an
automated assertion in the loaders; failures exit non-zero.$^*$ means pass for CRSP/iid while TAQ had only 29 names}
\label{tab:ledger}
\scriptsize
\begin{tabular}{@{}l l l@{}}
\toprule
Rule & Check & Status \\
\midrule
R0 & CRSP parquet exists, 92{,}681 rows & pass \\
R1 & iid parquet exists, 93{,}128 rows  & pass \\
R2 & CRSP and iid share calendar window & pass \\
R3 & PIT merge yields 30 names per day  & pass$^*$ \\
R4 & \texttt{ret} vs.\ \texttt{retx} divergence plausible (1.5\% of rows) & pass \\
R5 & Permno map covers DD/DWDP/DOW and UTX/RTX chains & pass \\
R6 & $\geq$30 earnings announcements per always-present ticker & pass \\
R7 & 35/35 TAQ aggregation unit tests pass & pass \\
R8 & 92{,}534 / 92{,}535 H1 ticker-days have correct bar count & pass \\
R9 & 5/5 sample-day VWAPs agree with CRSP within 5\% & pass (0.3\% max) \\
R10 & Backtest engine mirror invariant $\Pi^{\MOM}+\Pi^{\REV}=0$ & pass ($<10^{-15}$) \\
R11 & Turnover $\tau_t \in [0, 2]$ on every rebalance & pass \\
\bottomrule
\end{tabular}
\end{table}

\section{Newey--West lag sensitivity for H6}
\label{app:nw}

The headline H6 MOM result uses $q=6$ Newey--West lags to correct
for the 5-month forward-window overlap. Table~\ref{tab:nw-sens}
sweeps $q \in \{0, 2, 4, 6, 8, 12, 18\}$. The corrected $t$ is
stable at $1.35$--$1.53$ for every $q \geq 4$, versus naive $t$ of
$+2.43$. The ratio is close to the theoretical $1/\sqrt{1+2\cdot5/6}
\approx 0.61$ under full 5-month overlap. H6 fails Bonferroni at
every choice of $q > 0$, so the ``fallback'' label does not depend
on an arbitrary lag selection.

\begin{table}[ht]
\centering
\caption{Newey--West-corrected $t$-statistic on H6 MOM across lag
choices. $q{=}0$ is the naive SE estimate. Bonferroni threshold
$t^{*} \approx 2.87$.}
\label{tab:nw-sens}
\small
\input{icm/tables/h6_nw_sensitivity}
\end{table}

\section{Stationary-bootstrap block-length sensitivity (H2 REV)}
\label{app:boot}

The stationary bootstrap of Politis--Romano~\cite{pr1994} draws
geometric-length blocks from the P\&L series to resample with
dependence preserved. Block length $\ell$ is a bias-variance
trade-off: too short misses serial dependence, too long yields
fewer effectively-independent draws. H2 REV's daily P\&L is
near-i.i.d.\ (Figure~\ref{fig:acf} shows all 30 autocorrelations
within $\pm 0.04$), so the bootstrap CI should be stable across
$\ell$. Table~\ref{tab:boot-sens} confirms: the lower bound
ranges from $+1.64$ ($\ell{=}100$) to $+1.88$ ($\ell \leq 10$);
the CI width grows modestly from 1.33 to 1.87 as $\ell$ increases.
At no block length does the CI touch zero.

\begin{table}[ht]
\centering
\caption{Bootstrap 95\% CI for annualised Sharpe of H2 REV at
several block lengths. Point estimate $+2.53$; $B = 2{,}000$
replicates, $N = 2{,}431$ daily observations.}
\label{tab:boot-sens}
\small
\input{icm/tables/h2_bootstrap_sensitivity}
\end{table}

\section{Half-spread-sourced cost model for H2 REV}
\label{app:halfspread}

The report's headline charges a flat $1.5$~bps per side
(round-trip $3$~bps). Appendix~D cross-checks this by driving the
cost model from iid\_ms's actual quoted spread per (ticker, date):
on each rebalance, we pay half the portfolio-weighted quoted spread
at entry and the same at exit. Table~\ref{tab:halfspread} reports
four scenarios.

\begin{table}[ht]
\centering
\caption{H2 REV under four cost models. ``Avg.\ cost'' is the
portfolio's realised round-trip cost per day in basis points.
Half-spread is the $|w|$-weighted mean of iid\_ms
\texttt{QuotedSpread\_Percent\_tw} across the day's traded names.}
\label{tab:halfspread}
\scriptsize
\input{icm/tables/h2_halfspread_cost}
\end{table}

The empirical weighted-average half-spread on the 20-name H2
portfolio is $1.52$~bps per side --- essentially identical to the
flat $1.5$~bps assumption in the report. That is not a coincidence:
the flat assumption was calibrated to the DJ30 large-cap half-spread
in the first place. Adding a $0.5$~bps per-side impact floor
(``half-spread + 0.5 bps impact'') raises cost to $4.03$~bps
round-trip; net Sharpe drops from $+2.53$ to $+1.90$ but $t{=}+5.90$
still clears Bonferroni by more than $2\times$.

\section{Backtest-engine invariants}
\label{app:invariants}

Four assertions run after every backtest:

\begin{enumerate}[leftmargin=*]
  \item \textbf{Mirror invariant.}
    $\bigl|\sum_i (w^{\MOM}_{i,t} + w^{\REV}_{i,t})\bigr| < 10^{-12}$
    for every rebalance. The two strategies are pure negatives of
    each other; any non-zero sum flags a weight-assignment bug.
  \item \textbf{Gross-leverage invariant.}
    $0.99 \leq \sum_i |w^{\MOM}_{i,t}| \leq 1.01$ for every
    rebalance (the $\pm 1\%$ tolerance covers short days where
    tercile sizes become $9/11/9$).
  \item \textbf{Turnover bound.}
    $\tau_t \in [0, 2]$ on every rebalance. A turnover above 2 is
    impossible for a gross-one dollar-neutral portfolio; below 0
    is a sign convention bug.
  \item \textbf{Membership closure.} Every $(t, i)$ in the P\&L
    frame satisfies $i \in U_t$, enforced by inner-join against
    the PIT membership table. Any violation would indicate a
    stale-universe bug.
\end{enumerate}

These four cover the failure modes that cost strategies their
Sharpe in practice (weight-leakage, leverage drift, double-counted
turnover, out-of-universe positions). All four assert silently on
every run of the 16-step regeneration pipeline; our
findings are the output of a pipeline that has passed each one
every time.

\section{What could go wrong that didn't}
\label{app:gremlins}

A short list of issues that \emph{should} have broken the
pipeline but did not, curated from the attribution logs. The
usefulness of the list is less about the specific bugs and more
about the class of bug each could have produced.

\paragraph{WRDS pull missing a DJ30 member.} The original extract
was 39 tickers instead of 40 (Goldman Sachs absent). Caught by the
loader's ``40 ever-members'' assertion, not by a downstream
symptom. Mitigation was a supplementary PERMNO-specific pull,
concatenated into the loader. If the assertion had been
``$\geq$1 ticker'' the entire backtest would have run on a 29-name
universe and the tercile sizes would silently have been $10/9/10$.

\paragraph{Multi-security iid rows collapsing.} iid\_ms uses
\texttt{SYM\_ROOT} only, so JPM preferreds (JPMPRA, \ldots, JPMPRM)
and DD class-B share the common-stock ticker. Caught by the
uniqueness check $(date, ticker)$. Resolved by sorting descending
on \texttt{total\_dollar\_m} and keeping the top row --- the
common stock outranks preferreds by 1--3 orders of magnitude so
the pick is unambiguous.

\paragraph{TIME\_M format drift.} TAQ ships times as
\texttt{H:MM:SS.fffffffff} with \emph{variable-width} hours; a
naive fixed-slice parser fails on 4:00~AM pre-market prints.
Caught by a pytest regression case after the first
cluster-job crash.

\paragraph{UTX$\to$RTX ticker-label mismatch.} CRSP and iid
relabelled PERMNO~17830 as RTX on merger day, but S\&P DJI kept
UTX in the index through that day. Caught by the 30-names
assertion failing on exactly one date (2020-04-03). Resolved
by a one-cell loader patch; the underlying entity is the same
security.

\paragraph{IBES legacy ticker field.} IBES's \texttt{TICKER}
column uses pre-acquisition symbols (NIKE instead of NKE, VISA
instead of V, XON instead of XOM). A naive join on \texttt{TICKER}
would have dropped 8 of the 24 always-present DJ30 names from
the earnings-blackout robustness check. Caught by the IBES
loader's ``$\geq$30 announcements per always-present ticker''
assertion; resolved by joining on \texttt{OFTIC} (the stable
trading symbol).

\paragraph{Intraday cost accounting.} H2's default turnover under
the carry-through cost model would have been $\sim\!1.3$ per day
rather than the correct $2.0$, understating cost by roughly a
factor of $1.5$. Caught during a review of the round-trip flag
on H2 and H1; resolved by adding the explicit intraday
``full round-trip'' branch in Algorithm~\ref{alg:engine}. If
this had stayed wrong, H2 REV's reported net Sharpe would have
been $\approx +3.1$ instead of $+2.53$ --- too optimistic.

\section{Further work}
\label{app:further}

Three extensions that were out of scope for the first pass:

\begin{enumerate}[leftmargin=*]
  \item \emph{Short-sale constraints.} The strategy assumes costless
        shorting. A realistic model of borrow cost on DJ30 names
        ($\sim$25--50~bps annualised for hard-to-borrow, $\sim$5~bps
        for easy-to-borrow) would shift each short leg. DJ30 is
        predominantly easy-to-borrow so the impact is modest; a
        worst-case 50~bps annualised on the short leg reduces net
        Sharpe by $\sim 0.08$.
  \item \emph{Participation-scaled cost curve.} The half-spread
        cost model in Appendix~D treats cost as signal-independent.
        Realistic execution at size includes a participation-scaled
        impact term (Almgren--Chriss square-root law or similar).
        Calibrating an impact model on the DJ30 volume profile
        would refine the capacity discussion in \S\ref{sec:impl}.
  \item \emph{Machine-learning uplift.} The current signal is a
        single rank. A gradient-boosted-tree regressor on the
        same panel, with features including lagged returns,
        trailing volatility, and order-imbalance proxies, might
        extract additional edge. Purged walk-forward validation
        with a 5-day embargo would be the honest measurement. The
        null hypothesis is no uplift; we'd test it by Diebold--Mariano
        on the paired P\&L series against the naive-rank baseline.
\end{enumerate}

%======================================================================
\bibliographystyle{plain}
\begin{thebibliography}{99}

\bibitem{lehmann1990}
B.~N.~Lehmann.
\newblock Fads, martingales, and market efficiency.
\newblock {\em Quarterly Journal of Economics}, 105(1):1--28, 1990.

\bibitem{lo1990}
A.~W.~Lo and A.~C.~MacKinlay.
\newblock When are contrarian profits due to stock market overreaction?
\newblock {\em Review of Financial Studies}, 3(2):175--205, 1990.

\bibitem{jegadeesh1990}
N.~Jegadeesh.
\newblock Evidence of predictable behavior of security returns.
\newblock {\em Journal of Finance}, 45(3):881--898, 1990.

\bibitem{jt1993}
N.~Jegadeesh and S.~Titman.
\newblock Returns to buying winners and selling losers: implications
  for stock market efficiency.
\newblock {\em Journal of Finance}, 48(1):65--91, 1993.

\bibitem{gao2018}
L.~Gao, Y.~Han, S.~Z.~Li, and G.~Zhou.
\newblock Market intraday momentum.
\newblock {\em Journal of Financial Economics}, 129(2):394--414, 2018.

\bibitem{roll1984}
R.~Roll.
\newblock A simple implicit measure of the effective bid-ask spread
  in an efficient market.
\newblock {\em Journal of Finance}, 39(4):1127--1139, 1984.

\bibitem{lm1988}
A.~W.~Lo and A.~C.~MacKinlay.
\newblock Stock market prices do not follow random walks: evidence
  from a simple specification test.
\newblock {\em Review of Financial Studies}, 1(1):41--66, 1988.

\bibitem{nw1987}
W.~K.~Newey and K.~D.~West.
\newblock A simple, positive semi-definite, heteroskedasticity and
  autocorrelation consistent covariance matrix.
\newblock {\em Econometrica}, 55(3):703--708, 1987.

\bibitem{ff2015}
E.~F.~Fama and K.~R.~French.
\newblock A five-factor asset pricing model.
\newblock {\em Journal of Financial Economics}, 116(1):1--22, 2015.

\bibitem{carhart1997}
M.~M.~Carhart.
\newblock On persistence in mutual fund performance.
\newblock {\em Journal of Finance}, 52(1):57--82, 1997.

\bibitem{mg1999}
T.~J.~Moskowitz and M.~Grinblatt.
\newblock Do industries explain momentum?
\newblock {\em Journal of Finance}, 54(4):1249--1290, 1999.

\bibitem{dt1985}
W.~F.~M.~De~Bondt and R.~Thaler.
\newblock Does the stock market overreact?
\newblock {\em Journal of Finance}, 40(3):793--805, 1985.

\bibitem{nagel2012}
S.~Nagel.
\newblock Evaporating liquidity.
\newblock {\em Review of Financial Studies}, 25(7):2005--2039, 2012.

\bibitem{lo2004}
A.~W.~Lo.
\newblock The adaptive markets hypothesis.
\newblock {\em Journal of Portfolio Management}, 30(5):15--29, 2004.

\bibitem{pr1994}
D.~N.~Politis and J.~P.~Romano.
\newblock The stationary bootstrap.
\newblock {\em Journal of the American Statistical Association},
  89(428):1303--1313, 1994.

\bibitem{nysetaq}
NYSE.
\newblock Daily TAQ client specification, version 2.2a, 2022.

\end{thebibliography}

\end{document}